\chapter{Introduction and Inspiration}
\label{ch:introduction}

\section{Motivation}

The Cubilox Mk1 is the author's first large electronics project. Before it, only small Arduino
projects had been built. The idea for the cube goes back several years, to a video of the
\emph{Cubli}, a cube developed at the Institute for Dynamic Systems and Control of ETH Zurich that
can jump up and balance on its corner~\cite{gajamohan2012cubli,gajamohan2013cubli}. The goal was
to build a similar cube completely independently, and to learn the required electronics, control
engineering and mechanical design along the way.

While researching what such a cube needs, the author came across the self-balancing cube of
remrc~\cite{remrc_cube}. This project had a strong influence on the Cubilox Mk1: it uses the
same type of motor (Nidec~24H) and a similar general layout. Nevertheless, every part of the
Cubilox Mk1 was created by the author: all mechanical parts were designed from scratch, the
circuit was designed and soldered independently, and the firmware was written without using code
from other cube projects.

Because it is a first project of this size, not every solution is best practice, and some
decisions would be made differently today. These points are discussed openly in
\cref{ch:limitations,ch:outlook}. The cube works, and the knowledge gained here is the basis for
future projects.

\section{Scientific background: the Cubli}

The Cubli is a \SI{15}{cm} cube with reaction wheels mounted on three of its faces. By applying
controlled torques to the wheels, it balances on an edge or a corner; by spinning the wheels up and
braking them abruptly, it can even jump up from a resting position~\cite{gajamohan2012cubli}. In
the second publication, the multi-body dynamics of the Cubli are derived, its parameters are
identified, and corner balancing with a linear state feedback controller is demonstrated
experimentally~\cite{gajamohan2013cubli}. The Cubli can therefore be seen as a three-dimensional
version of the classical inverted pendulum, actuated purely by internal reaction wheels.

\section{Goals and scope of the Cubilox Mk1}

The goals of the project were:
\begin{itemize}
  \item a cube that balances on an edge using one reaction wheel,
  \item a cube that balances on a vertex using all three reaction wheels,
  \item a design that can be printed on a standard consumer 3D printer and built from easily
        available parts,
  \item self-written firmware that can be tuned and calibrated wirelessly over Bluetooth.
\end{itemize}
Jumping up from a resting position, as the Cubli does, was not a goal of the Mk1 (see
\cref{ch:outlook}).

\Cref{fig:render,fig:vertex_photo} show the CAD model and the finished cube balancing on its
vertex.

\begin{figure}[htbp]
  \centering
  \begin{subfigure}[b]{0.48\textwidth}
    \centering
    \includegraphics[width=\textwidth]{cubilox_mk1_render.png}
    \caption{Rendering of the CAD model.}
    \label{fig:render}
  \end{subfigure}\hfill
  \begin{subfigure}[b]{0.40\textwidth}
    \centering
    \includegraphics[width=\textwidth]{vertex_balance.jpg}
    \caption{The finished cube balancing on its vertex.}
    \label{fig:vertex_photo}
  \end{subfigure}
  \caption{The Cubilox Mk1: CAD model and built prototype.}
  \label{fig:overview}
\end{figure}

\section{Structure of this document}

\Cref{ch:principle} explains the physics of reaction wheel balancing.
\Cref{ch:mechanical,ch:electronics} describe the mechanical design and the electronics.
\Cref{ch:firmware} explains the firmware and the control laws in detail, followed by a getting
started guide (\cref{ch:getting_started}) and a tuning guide (\cref{ch:tuning}).
\Cref{ch:assembly} gives assembly instructions. \Cref{ch:results} presents the results and the
development history, \cref{ch:limitations} the limitations, and \cref{ch:outlook} the ideas for a
Mk2. The bill of materials and the third-party software are listed in the appendix.
